\documentclass[french,11pt]{book}
\usepackage{teach}
\usepackage[french]{babel}
\usepackage{amsmath,amssymb}
\usepackage{array}
\newcommand{\R}{\mathbb{R}}
\begin{document}
\begin{center}
{\Large \textbf{Exercices -- Série statistiques à deux variables}}\\[2mm]
Terminale technologique
\end{center}

\section*{Exercices}
\begin{enumerate}
  \item On donne les couples $(1;12)$, $(2;15)$, $(3;17)$, $(4;20)$, $(5;24)$. Représenter le nuage de points et commenter la tendance.
  \item À la calculatrice, déterminer une équation de la droite d'ajustement affine pour les données précédentes.
  \item Utiliser l'ajustement pour estimer la valeur de $y$ lorsque $x=8$.
  \item Une série donne un coefficient de corrélation $r=0,94$. Interpréter la qualité de l'ajustement affine.
\end{enumerate}

\end{document}
